\section{Cancellation in one time dimension}
\label{sec:cancellation}

The structural argument has produced finite sets of orthogonal lines whose
motion is described by an automorphism of an abstract abelian group.
We now prove the cancellation statement needed to rule out a nonzero
correlation with the zero-entropy witness. The probability law in this section
need not be invariant under any transformation of the space carrying the
lines. Only its marginal on the original process is used.

Let $(X,\mu_X,T_X)$ be a mixing system, let $(Y,\eta)$ be a probability
space, and let $x:Y\to X$ have law $\mu_X$. Suppose that a countable abelian
group $L$ is represented by measurable functions $\chi_l:Y\to\C$ such that
\begin{equation}
 \label{eq:projective-characters}
 |\chi_l|=1,\qquad \chi_0=1,\qquad
 \chi_l\chi_{l'}\in\T\chi_{l+l'},\qquad
 \int\chi_l\dd\eta=0\quad(l\ne0).
\end{equation}
Here $\T\chi$ denotes all constant multiples of $\chi$ by scalars of
modulus one; in particular, the scalar in the product relation is independent
of $y$. These assumptions imply that distinct $\chi_l$ are orthogonal.
Let $G:L\to L$ be a group automorphism.

\begin{proposition}[Orbit cancellation]
\label{prop:orbit-cancellation}
Assume that $X$ is a mixing two-sided finite-alphabet stationary process whose
alphabet partition has entropy rate zero. For every centered finite-valued
one-site function $f$ on $X$, every probability coupling $x:Y\to X$, every
family satisfying \eqref{eq:projective-characters}, every automorphism $G$,
and every $l\in L$, one has
\begin{equation}
 \label{eq:orbit-cancellation}
 \lim_{N\to\infty}\frac1N\sum_{r=0}^{N-1}
 \left|\int f(T_X^r x)\overline{\chi_{G^r l}}\dd\eta\right|=0.
\end{equation}
\end{proposition}

The proof separates three possibilities for a label orbit. A dissociated
orbit is handled by entropy. A finite relation either produces dissociated
subsequences on arithmetic progressions, or makes the orbit polynomial on
such progressions. Mixing then handles the polynomial case.

\subsection{Entropy against dissociated labels}

\begin{definition}
A sequence $(l_m)_{m\ge0}$ in an abelian group is \emph{dissociated} if
\[
 \sum_{m\in A}\varepsilon_m l_m\ne0
\]
for every finite nonempty set $A\subset\N\cup\{0\}$ and every choice
$\varepsilon_m\in\{-1,1\}$.
\end{definition}

This definition concerns distinct indices. Thus it excludes zero labels and
repetitions, and also applies to groups with torsion.

The proof begins with the exponential-moment argument for dissociated
characters; see \cite[Section~2.2, Proposition~17]{Green2002SetAddition}.
We give the argument for the projective characters above and then use
the sublinear word-entropy hypothesis.

\begin{lemma}[Entropy cancellation]
\label{lem:entropy-cancellation}
Suppose \eqref{eq:projective-characters} holds and $(l_m)$ is dissociated.
Let $(F_m)_{m\ge0}$ be complex random variables on $(Y,\eta)$ taking values
in a fixed finite subset of the closed unit disk. If the Shannon entropy
$H_n$ of $(F_0,\ldots,F_{n-1})$ satisfies $H_n=o(n)$, then
\[
 \frac1n\sum_{m<n}\left|\int F_m\overline{\chi_{l_m}}\dd\eta\right|
 \longrightarrow0.
\]
No independence between the variables $F_m$ and the functions $\chi_{l_m}$
is assumed.
\end{lemma}

\begin{proof}
For deterministic complex numbers $w_m$ with $|w_m|\le1$ and
$0<\theta\le1$, the elementary inequality
$e^t\le1+t+2t^2$ for $|t|\le1$ gives
\[
 \exp\bigl(\theta\Re(w_m\overline{\chi_{l_m}})\bigr)
 \le 1+\theta\Re(w_m\overline{\chi_{l_m}})+2\theta^2.
\]
The right side is nonnegative. Multiply these inequalities over $m<n$ and
integrate. On expanding the product, a nonconstant term contains one factor
$\chi_{l_m}$ or $\overline{\chi_{l_m}}$ at each of a nonempty set of
distinct indices. By \eqref{eq:projective-characters}, it is a constant
multiple of the representative of a signed sum of those labels.
Dissociation makes that label nonzero, so its integral vanishes. Hence
\begin{equation}
 \label{eq:dissociated-exponential}
 \int\exp\left(\theta\sum_{m<n}
       \Re(w_m\overline{\chi_{l_m}})\right)\dd\eta
 \le(1+2\theta^2)^n\le e^{2\theta^2 n}.
\end{equation}
For $0<\varepsilon\le1$, choose $\theta=\varepsilon/4$ and apply
Markov's inequality. Uniformly over the deterministic coefficients,
\begin{equation}
 \label{eq:dissociated-tail}
 \eta\left\{\frac1n\sum_{m<n}
       \Re(w_m\overline{\chi_{l_m}})>\varepsilon\right\}
 \le e^{-\varepsilon^2 n/8}.
\end{equation}

We next allow the coefficients to be the random word $(F_m)$. Put
\[
 a_n=\sqrt{H_n/n}+n^{-1/2}.
\]
Call a word likely if its probability is at least $e^{-a_n n}$. There are
at most $e^{a_n n}$ likely words. For the random word, the mean of minus
the logarithm of its probability is $H_n$, so the total probability of
the other words is at most $H_n/(a_n n)$.
Both $a_n$ and this last bound tend to zero.

Fix arbitrary deterministic phases $\omega_m\in\T$. On the event that
$(F_m)_{m<n}$ equals a particular likely word $(v_m)$, the sum
\[
 Y_n=\frac1n\sum_{m<n}
       \Re(\omega_m F_m\overline{\chi_{l_m}})
\]
is the sum in \eqref{eq:dissociated-tail} with $w_m=\omega_m v_m$.
We do not condition the character law on that word: the event in question
is simply a subset of the corresponding event for deterministic
coefficients. A union bound therefore gives
\begin{equation}
 \label{eq:likely-word-bound}
 \eta\{Y_n>\varepsilon\}
 \le \frac{H_n}{a_n n}
       +\exp\left(a_n n-\frac{\varepsilon^2 n}{8}\right)
 \longrightarrow0.
\end{equation}
This explains why arbitrary coupling is allowed.

Finally choose each $\omega_m$ to make
$\omega_m\int F_m\overline{\chi_{l_m}}\dd\eta$ nonnegative real,
choosing any phase if that integral is zero. These choices are
deterministic. Since $Y_n\le1$, we have
\[
 0\le\frac1n\sum_{m<n}
       \left|\int F_m\overline{\chi_{l_m}}\dd\eta\right|
   =\int Y_n\dd\eta
   \le\varepsilon+\eta\{Y_n>\varepsilon\}.
\]
First let $n$ tend to infinity and then let $\varepsilon$ tend to zero.
\end{proof}

For the process in Proposition~\ref{prop:orbit-cancellation}, fix integers
$p\ge1$ and $0\le j<p$. The word
$(f(T_X^{pm+j}x))_{m<n}$ is a function of a consecutive alphabet window
of length at most $pn+p$. Stationarity and zero entropy give entropy
$o(n)$ for this word. After scaling $f$, Lemma~\ref{lem:entropy-cancellation}
therefore proves \eqref{eq:orbit-cancellation} on any arithmetic progression
on which its label sequence is dissociated.

\subsection{The algebra of a single orbit}

The finite-generation and spectral steps below have predecessors in
Rokhlin's compact-group argument \cite[Sections~2.3 and~3.2]{Rokhlin1949}.
Here the labels may have torsion, and the case in which a power acts
unipotently is retained for polynomial cancellation.

\begin{lemma}[Orbit alternatives]
\label{lem:orbit-alternatives}
Let $G$ be an automorphism of an abelian group $L$ and $l\in L$.
At least one of the following holds:
\begin{enumerate}
\item The sequence $(G^r l)_{r\ge0}$ is dissociated.
\item There is $p\ge1$ such that $(G^{pm+j}l)_{m\ge0}$ is dissociated
      for every $0\le j<p$.
\item There are $p,d\ge1$ such that, on the subgroup
      $K=\langle G^r l:r\in\Z\rangle$, the endomorphism
      $N=G^p-\id$ satisfies $N^d=0$. In this case
      \begin{equation}
       \label{eq:binomial-orbit}
       G^{pm+j}l
       =\sum_{a=0}^{d-1}\binom ma N^aG^j l
       \qquad(m\ge0,\ 0\le j<p).
      \end{equation}
\end{enumerate}
\end{lemma}

\begin{proof}
Suppose the first alternative fails. Remove any initial zero coefficients
from a nonempty signed relation and apply a power of $G^{-1}$. We obtain
\[
 b_0l+b_1Gl+\cdots+b_qG^q l=0,
 \qquad b_0,b_q\in\{-1,1\},\quad b_a\in\{-1,0,1\}.
\]
If $q=0$, then $l=0$ and the third alternative holds with $p=d=1$.
For $q\ge1$, the unit coefficient at the upper end expresses $G^ql$
as an integral combination of the preceding $q$ labels. The coefficient
at the lower end does the same in the other direction. Applying the
relation to all translates shows that $K$ is generated by
$l,Gl,\ldots,G^{q-1}l$. In particular, its torsion subgroup $K_{\rm tor}$
is finite and its quotient $\overline K$ is a finite-rank lattice.

The induced map $\overline G$ is a lattice automorphism. If it has an
eigenvalue $\lambda$ with $|\lambda|>1$, choose a nonzero complex left
eigenfunctional $\varphi$ on $\overline K\otimes_\Z\C$. The two-sided
orbit of $\overline l$ spans this vector space. Thus
$\varphi(\overline l)\ne0$, since otherwise
$\varphi(\overline G^r\overline l)=\lambda^r
\varphi(\overline l)=0$ for every integer $r$.
Choose $p$ with $R=|\lambda|^p>2$.
On each residue class $j$, the absolute values of the projected labels
are $|\lambda^j\varphi(\overline l)|R^m$. In any finite signed relation
on distinct indices, the term with largest index strictly dominates the
sum of all earlier terms, because $\sum_{a<m}R^a<R^m$.
There can be no such relation in $K$. This proves the second alternative.

It remains to consider the case with no expanding eigenvalue. If the
lattice has positive rank, $\det\overline G=\pm1$ implies that every
eigenvalue has modulus one. We use the bounded-coefficient argument in
Kronecker's root-of-unity criterion \cite[Theorem~I]{Kronecker1857}; see also
\cite[Theorem~1 and its proof]{Damianou2015}.
The characteristic polynomial of
$\overline G^n$ has integer coefficients, each bounded independently of
$n$: its coefficients are elementary symmetric functions of a fixed
number of numbers of modulus one. Only finitely many such polynomials
exist. For each eigenvalue $\lambda$, all the powers $\lambda^n$ are
roots of these finitely many polynomials. Hence two powers agree and
$\lambda$ is a root of unity.

Choose $p$ such that $\overline G^p$ is unipotent and $G^p$ is the
identity on $K_{\rm tor}$. The latter is possible because the automorphism
group of a finite group is finite. This argument includes the rank-zero
case. For some $d_0\ge1$, the map $N=G^p-\id$ satisfies
$N^{d_0}K\subseteq K_{\rm tor}$, and $NK_{\rm tor}=0$.
Consequently $N^{d_0+1}K=0$.
The additional application of $N$ accounts for possible maps from a free
part into torsion; an invariant splitting of $K$ has not been assumed.
Taking $d=d_0+1$ and expanding $(\id+N)^m$ proves
\eqref{eq:binomial-orbit}.
\end{proof}

\subsection{Polynomial labels and mixing}

A group-valued polynomial of degree at most $d$ on the nonnegative
integers will mean a sequence
\[
 P(m)=\sum_{a=0}^d\binom ma l_a,\qquad l_a\in L.
\]
For each fixed integer $h\ge1$, the difference $P(m+h)-P(m)$ has degree
at most $d-1$. Indeed, Vandermonde's identity expands
$\binom{m+h}{a}-\binom ma$ as an integral combination of
$\binom m0,\ldots,\binom m{a-1}$. This observation is valid in every
abelian group, including groups with torsion.

We use Hilbert-space van der Corput differencing and induction on
polynomial degree; compare \cite[Theorems~1.4--1.5 and p.~339]{Bergelson1987}.
The argument below works in the present coupling, whose law need not
be invariant.

\begin{lemma}[Polynomial cancellation]
\label{lem:polynomial-cancellation}
Let $(X,\mu_X,T_X)$ be mixing, let $x:Y\to X$ have law $\mu_X$, and
assume \eqref{eq:projective-characters}. For every bounded centered
function $F$ on $X$, every $p\ge1$, and every group-valued polynomial
$P(m)$, one has
\[
 \frac1n\sum_{m<n}
 \left|\int F(T_X^{pm}x)\overline{\chi_{P(m)}}\dd\eta\right|
 \longrightarrow0.
\]
Here neither zero entropy nor invariance of $\eta$ is required.
\end{lemma}

\begin{proof}
We induct on the degree. If $P(m)=l_0$ is constant, condition
$\overline{\chi_{l_0}}$ onto $x$. Its conditional expectation is a
bounded measurable function $g$ on $X$, and
\[
 \int F(T_X^{pm}x)\overline{\chi_{l_0}}\dd\eta
   =\int F(T_X^{pm}x)g(x)\dd\mu_X(x)\longrightarrow0
\]
by mixing and centering. Mixing for bounded functions follows from the
set formulation by simple-function approximation. This proves the base
case, including $l_0=0$.

For the inductive step, choose deterministic phases $\omega_m$ aligning
the integrals in the assertion. In the fixed Hilbert space $L^2(\eta)$
put
\[
 v_m=\omega_m F(T_X^{pm}x)\overline{\chi_{P(m)}}.
\]
Their norms are bounded by $\norm F_\infty$. At a fixed positive lag $h$,
the absolute value of $\ip{v_m}{v_{m+h}}$ equals
\[
 \left|\int F_h(T_X^{pm}x)
       \overline{\chi_{P(m+h)-P(m)}}\dd\eta\right|,
 \qquad F_h=(F\circ T_X^{ph})\overline F.
\]
The constant phases from alignment and from
\eqref{eq:projective-characters} disappear inside this absolute value.
Apply the induction hypothesis to the centered function
$F_h-\int F_h\dd\mu_X$ and the polynomial of lower degree.
The remaining constant contribution has absolute value at most
$|\int F_h\dd\mu_X|$. Therefore
\begin{equation}
 \label{eq:lag-correlation-bound}
 \limsup_{n\to\infty}\frac1n\sum_{m<n}
       |\ip{v_m}{v_{m+h}}|
 \le\left|\int F_h\dd\mu_X\right|.
\end{equation}

The following averaging step is the Hilbert-space van der Corput
argument; see \cite[Theorem~4.1]{BergelsonKnutson2008} for a real
Hilbert-space formulation, using the real part of the inner product in
the complex case. For completeness, fix $H\ge1$. Replacing each $v_m$
in its mean by
$H^{-1}\sum_{a=0}^{H-1}v_{m+a}$ changes the mean by a vector of norm
at most $2H\norm F_\infty/n$. Convexity of the squared norm bounds the
squared norm of this new mean by the mean of the squared norms of its
summands. Expanding those squares, grouping terms by the lag, and using
\eqref{eq:lag-correlation-bound} gives
\[
 \limsup_{n\to\infty}
 \left\lVert\frac1n\sum_{m<n}v_m\right\rVert^2
 \le\frac{\norm F_\infty^2}{H}
    +\frac2H\sum_{h=1}^{H-1}
       \left(1-\frac hH\right)
       \left|\int F_h\dd\mu_X\right|.
\]
The harmless shift of the averaging range by at most $H$ indices does
not change any of these limiting bounds. Mixing and centering imply
$\int F_h\dd\mu_X\to0$ as $h\to\infty$. Letting $H\to\infty$ shows
that the Hilbert-space means tend to zero.
Testing them against the constant function $1$ yields exactly the mean
of the aligned absolute integrals, proving the induction.
\end{proof}

\begin{proof}[Proof of Proposition~\ref{prop:orbit-cancellation}]
Apply Lemma~\ref{lem:orbit-alternatives}. In the first alternative,
Lemma~\ref{lem:entropy-cancellation} applies directly. In the second,
apply it on each of the finitely many residue classes modulo $p$, using
the progression-word entropy estimate after that lemma. In the third,
\eqref{eq:binomial-orbit} gives a polynomial on every residue class.
Apply Lemma~\ref{lem:polynomial-cancellation} there with
$F=f\circ T_X^j$. This function is centered by invariance of $\mu_X$.
In each case the finitely many progression averages tend to zero and
combine to give the full average in \eqref{eq:orbit-cancellation}.
\end{proof}
